\section{Introduction}

\subsection{Background and Motivation}

In recent years, local AI applications on consumer-grade devices have grown rapidly. Users increasingly expect multiple AI services to remain available on the same personal device for privacy, responsiveness, and offline availability. Typical scenarios include combining assistant-style interaction with intermittent burst requests such as image generation or document processing. However, such devices remain severely resource-constrained: mainstream laptops are often equipped with only 8GB VRAM, which in our setting is not enough to keep both a resident LLM runtime and a bursty GPU-intensive generation runtime resident at the same time. The key difficulty is not simply that the two models are different, but that they impose different service obligations: the resident LLM service should remain reachable during contention, whereas the burst task may wait but should obtain exclusive accelerator access once admitted. This mismatch between limited dedicated VRAM and asymmetric service obligations motivates a lightweight scheduler for local AI runtimes.

\subsection{Core Challenges}

\textbf{Hard Memory Constraints}. In our tested configuration, 8GB VRAM is insufficient to keep both the quantized LLM runtime and the image model resident at the same time, making traditional GPU isolation approaches impractical. We measure the full serialized migration state, the pre-allocated GPU KV buffer, and the whole-card co-located peak separately; the \texttt{save\_state()} payload is never used as a proxy for GPU KV VRAM. When a GPU-intensive burst task arrives, the VRAM occupied by the resident task must be released or downgraded quickly to avoid CUDA OOM.

\textbf{Asymmetric Service Requirements}. The two workload roles in our target setting behave differently at the service level, as shown in Table~\ref{tab:task_characteristics}:

\begin{table}[htbp]
\caption{Task Characteristics Comparison}
\label{tab:task_characteristics}
\centering
\small
\begin{adjustbox}{max width=\columnwidth}
\begin{tabular}{@{}lcccc@{}}
\toprule
Task Type & Resource & Latency & Queueable & Semantics \\
 & Requirement & Sensitivity &  &  \\
\midrule
LLM Chat & GPU/CPU & High (streaming) & Low & Resident \\
Image Gen & GPU-required & Medium (e2e) & High & One-shot \\
\bottomrule
\end{tabular}
\end{adjustbox}
\end{table}

The important point is not to build a broad taxonomy of local AI tasks, but to isolate the concrete asymmetry that drives scheduling decisions in this paper: one workload should stay reachable under contention, while the other can queue but needs exclusive GPU ownership to complete efficiently.

\subsection{Limitations of Existing Solutions}

Existing solutions mainly target two different design centers: (1) distributed scheduling frameworks (e.g., Ray, Kubernetes), which address cluster-level placement and orchestration but are generally mismatched to single-node consumer devices; and (2) single-model inference optimization systems (e.g., vLLM, FlexGen), which improve throughput or offloading efficiency for one model at a time but do not directly address service-level arbitration between multiple local services sharing the same constrained GPU.

The scheduling problem for mixed local AI services on consumer hardware remains less studied than data-center scheduling and single-model serving. Existing solutions do not directly target contention between a continuously available, migratable service and a bursty GPU-exclusive job on 8GB-class devices, nor do they provide much support for maintaining availability during temporary degradation. EdgeSched targets this gap by focusing on the control consequence of this asymmetry: releasing GPU resources quickly on task arrival is a different problem from restoring high-performance service after the burst task finishes.

More generally, we frame this setting as a \emph{two-role contention model} rather than as a hard-coded pair of applications. One role is a \emph{resident service}: latency-sensitive, long-lived, and allowed to degrade temporarily instead of becoming unavailable. The other is a \emph{burst service}: queueable, shorter lived, and admitted only when it can obtain exclusive GPU ownership. LLM serving and image generation are the concrete instantiation used in this paper, but the scheduling argument is expressed at the level of these two service roles.

\subsection{Our Contributions}

This paper presents EdgeSched, a local scheduling framework for mixed AI workloads on consumer-grade devices, with core contributions including:

\begin{enumerate}
\item \textbf{A Two-Role Contention Abstraction for Local Runtimes}. We characterize the target setting using a minimal service-role asymmetry: a resident interactive service should remain reachable under contention, whereas a burst task is queueable but requires exclusive GPU access once admitted. We use this abstraction to derive scheduling decisions in the studied setting, without claiming a universal ontology for all local AI workloads.

\item \textbf{Asymmetric Resource Migration Policy}. We show that, in this setting, symmetric migration does not jointly optimize both directions well. Downlink prioritizes fast VRAM release for memory safety, which favors immediate interruption and aggressive GPU state teardown; uplink prioritizes minimizing user-visible interruption, which favors background preloading and delayed cut-over. Because these objectives pull the migration procedure in opposite directions, a single symmetric policy faces a strong trade-off between the two paths. EdgeSched therefore uses asymmetric policies as a constraint-driven design choice rather than as a cosmetic optimization.

\item \textbf{Lifecycle-Aware Runtime Control}. We combine priority queues, backpressure control, cancellation semantics, and crash recovery (about 44 ms for semantic recovery in our measurements) to provide predictable scheduling behavior under mixed workloads and failures.

\item \textbf{Empirical Validation Against Practical Alternatives}. We evaluate the framework on an RTX 5070 Laptop (8GB VRAM) and compare it against restart-based, generic memory-management, and migration-based alternatives, showing that these simpler designs do not simultaneously provide continuity, low-interruption recovery, and bounded restoration cost in our deployment setting.
\end{enumerate}

EdgeSched addresses an underexplored scheduling problem in local AI runtimes, instantiated here with LLM serving and image generation as a concrete high-conflict workload pair under dedicated-VRAM constraints.

\textbf{Reproducibility}. The source code and experimental scripts are included in the supplementary material.
